\section{Quenched representation of an imaginary-time propagator\label{sec:quenched}}

The Hamiltonian of mean force is defined through a Euclidean object,
$\bar\rho_Q(\beta)=\Tr_X e^{-\beta H_{\mathrm{tot}}}$, whose natural representation
is an imaginary-time path integral~\cite{feynmanQuantumMechanicsPath1965, kleinertPathIntegralsQuantum2009}. While this representation makes exact treatment of the bath possible, it comes at a conceptual cost: the path
integral trades operator structure for functional structure. In particular,
questions about the \emph{operator algebra} underlying the reduced equilibrium
state are obscured once the dynamics is expressed solely in terms of paths and
kernels.

If one wishes to understand when a local or otherwise restricted Hamiltonian of
mean force can exist, it is therefore useful to reconnect the path-integral
description to an operator-based notion of imaginary-time propagation. The key
observation is that the reduced equilibrium operator may be regarded as an
imaginary-time evolution over an interval of length $\beta$, but with a generator
that need not be time independent. Allowing for explicit $\tau$-dependence in
the generator provides a natural bridge between the nonlocal Euclidean influence
functional and a time-ordered exponential acting on the system Hilbert space.

Motivated by this, we introduce a class of imaginary-time propagators generated
by a $\tau$-dependent Hamiltonian on the system alone. These propagators will be
constructed for fixed external fields and only subsequently related to the
microscopic model. We refer to this construction as \emph{quenched} to emphasise
that the external field is held fixed during the imaginary-time evolution,
rather than being averaged over inside the exponential. In this sense, the
imaginary-time generator encodes an effective, path-dependent temperature
profile whose statistics are imposed only after the operator has been defined.

This quenched representation serves two purposes. First, it provides a
well-defined operator object whose path-integral representation can be derived
by standard means, allowing a precise comparison with the Euclidean influence
functional. Second, it isolates the algebraic content of the problem: all
nontrivial operator structure is carried by the imaginary-time–dependent
generator, while nonlocality and temperature dependence enter only through the
external field. In later sections we show how, for Gaussian environments, the
microscopic bath generates a specific quenched ensemble of such propagators,
and how the commuting case collapses this structure to a closed-form Hamiltonian
of mean force.

Let $H_Q$ be a system Hamiltonian and let $f$ be an arbitrary system operator.
For a prescribed c-number field $\xi(\tau)$ on $\tau\in[0,\beta]$ we define the
\emph{quenched} imaginary-time propagator
\begin{equation}
    \tilde\rho_Q(\beta;\xi)
    :=
    \mathcal T_\tau
    \exp\!\left[
        -\int_0^\beta d\tau\,
        \big(H_Q+\xi(\tau)\,f\big)
    \right].
    \label{eq:quenched_operator_def}
\end{equation}
The term \emph{quenched} refers only to the status of $\xi(\tau)$ as an external
field held fixed while constructing $\tilde\rho_Q(\beta;\xi)$.

\subsection{Generic phase-space path integral}

We derive a path-integral representation for the kernel
\begin{equation}
    \tilde\rho_Q(q_f,q_i;\beta|\xi)
    := \langle q_f|\tilde\rho_Q(\beta;\xi)|q_i\rangle,
\end{equation}
without assuming that $f$ is diagonal in the $|q\rangle$ basis.

Discretise $[0,\beta]$ into $N$ slices of width $\epsilon=\beta/N$, with
$\tau_n=n\epsilon$ and $\xi_n:=\xi(\tau_n)$. Writing
\begin{equation}
    \tilde\rho_Q(\beta;\xi)
    =
    \lim_{N\to\infty}
    \prod_{n=0}^{N-1}
    \exp\!\left[-\epsilon\big(H_Q+\xi_n f\big)\right],
\end{equation}
insert $N-1$ resolutions of the identity $\mathbf 1=\int dq\,|q\rangle\langle q|$ to obtain
\begin{equation}
\begin{split}
    \tilde\rho_Q(q_f,q_i;\beta|\xi)
    &=
    \lim_{N\to\infty}
    \int \prod_{n=1}^{N-1} dq_n \\
    &\quad \times \prod_{n=0}^{N-1}
    \langle q_{n+1}|e^{-\epsilon(H_Q+\xi_n f)}|q_n\rangle,
\end{split}
    \label{eq:quenched_trotter_q}
\end{equation}
with $q_0=q_i$ and $q_N=q_f$.

To obtain a generic expression for the short-time kernel, insert a momentum resolution
$\mathbf 1=\int \frac{dp}{2\pi\hbar}\,|p\rangle\langle p|$ between $|q_n\rangle$ and
$|q_{n+1}\rangle$, and use a mid-point (Weyl) prescription. For sufficiently small
$\epsilon$ one has the operator identity
\begin{equation}
\begin{split}
    \langle q'|e^{-\epsilon \hat A}|q\rangle
    &=
    \int\frac{dp}{2\pi\hbar}\,
    \exp\!\left[\frac{i}{\hbar}p(q'-q)\right]
    \\
    &\quad\times
    \exp\!\left[-\epsilon\,A_W\!\left(\frac{q'+q}{2},p\right)\right]
    +\mathcal O(\epsilon^2),
\end{split}
    \label{eq:short_time_weyl}
\end{equation}
where $A_W(q,p)$ denotes the Weyl symbol of the operator $\hat A$, corresponding to the symmetric ordering of operators in the underlying quantum theory~\cite{schulmanTechniquesApplicationsPath1981, berezinMethodSecondQuantization1966}. The mid-point evaluation of the coordinate argument, $(q'+q)/2$, is crucial here; other discretization prescriptions lead to different operator orderings in the continuum limit (e.g., pre-point evaluation yields standard ordering)~\cite{groscheHandbookFeynmanPath1998, kleinertPathIntegralsQuantum2009}. While the choice of ordering is formally a matter of convention, the Weyl prescription is particularly convenient for preserving symplectic symmetries in phase space.

\begin{widetext}
Applying \eqref{eq:short_time_weyl} with $\hat A = H_Q+\xi_n f$ gives
\begin{equation}
    \langle q_{n+1}|e^{-\epsilon(H_Q+\xi_n f)}|q_n\rangle
    = \int\frac{dp_n}{2\pi\hbar} \exp\!\left[\frac{i}{\hbar}p_n(q_{n+1}-q_n)\right] \exp\!\bigg[ -\epsilon \big(H_Q+\xi_n f\big)_W \left(\frac{q_{n+1}+q_n}{2},p_n\right) \bigg] +\mathcal O(\epsilon^2).
    \label{eq:short_time_quenched}
\end{equation}

Substituting \eqref{eq:short_time_quenched} into \eqref{eq:quenched_trotter_q} yields the discrete phase-space path integral
\begin{equation}
    \tilde\rho_Q(q_f,q_i;\beta|\xi) = \lim_{N\to\infty} \int \prod_{n=1}^{N-1} dq_n \int \prod_{n=0}^{N-1}\frac{dp_n}{2\pi\hbar} \exp\!\Bigg[ \frac{i}{\hbar}\sum_{n=0}^{N-1} p_n(q_{n+1}-q_n) -\epsilon\sum_{n=0}^{N-1} \big(H_Q+\xi_n f\big)_W\!\left(\frac{q_{n+1}+q_n}{2},p_n\right) \Bigg],
    \label{eq:quenched_discrete_phase_space}
\end{equation}
with fixed endpoints $q_0=q_i$, $q_N=q_f$. In the continuum limit this becomes
\begin{equation}
    \tilde\rho_Q(q_f,q_i;\beta|\xi) = \int_{q(0)=q_i}^{q(\beta)=q_f}\!\!\mathcal D q \int \mathcal D p \exp\!\left[ \frac{i}{\hbar}\int_0^\beta d\tau\, p(\tau)\,\dot q(\tau) -\int_0^\beta d\tau \big(H_Q+\xi(\tau) f\big)_W\!\big(q(\tau),p(\tau)\big) \right],
    \label{eq:quenched_phase_space_PI}
\end{equation}
\end{widetext}
where the subscript $W$ indicates the chosen operator-to-symbol map (here Weyl/mid-point).

Equation \eqref{eq:quenched_phase_space_PI} is the generic path-integral representation
of the quenched $\tau$-ordered propagator \eqref{eq:quenched_operator_def}. No assumption
has been made about diagonal structure of $f$; all operator ordering information is
carried by the fixed discretisation prescription defining the symbol
$\big(H_Q+\xi(\tau)f\big)_W$.
